\documentclass[11pt]{article}

\usepackage[margin=1in]{geometry}
\usepackage{amsmath}
\usepackage{titlesec}
\usepackage{enumitem}
\usepackage{listings}
\usepackage{xcolor}
\usepackage{booktabs}
\usepackage{hyperref}

\hypersetup{colorlinks=true, linkcolor=blue, urlcolor=blue}

\definecolor{codebg}{rgb}{0.96,0.96,0.96}
\lstset{
  basicstyle=\ttfamily\small,
  backgroundcolor=\color{codebg},
  breaklines=true,
  frame=single,
  columns=fullflexible,
  keepspaces=true,
}
\newcommand{\code}[1]{\texttt{#1}}
\newcommand{\ui}[1]{\textbf{#1}}

\titleformat{\section}{\large\bfseries}{\thesection}{1em}{}
\titleformat{\subsection}{\normalsize\bfseries}{\thesubsection}{1em}{}
\titleformat{\subsubsection}{\small\bfseries}{\thesubsubsection}{1em}{}

\title{Density-Based Traffic Signal Control \\ \large User Guide --- Navigating and Using the Application}
\author{Operator manual}
\date{2026-06-15}

\begin{document}
\maketitle

\tableofcontents
\bigskip

\section{Introduction}
This guide explains how to operate the desktop control panel after it has been
installed (see the Installation Guide). The application detects vehicles in up
to four camera feeds, decides which of the two traffic lights (the
\emph{vertical} lane and the \emph{horizontal} lane) should be green, and
publishes those decisions over MQTT. It also lets emergency vehicles
(ambulances, firetrucks) preempt the signal.

\subsection{Launching}
From the repository root, with the virtual environment active:
\begin{lstlisting}[language=bash]
source env/bin/activate
python gui/main.py
\end{lstlisting}
The main window opens with the title \ui{Control Panel}. All settings are
remembered between runs, so on later launches the app reopens with your last
cameras, layout, and options.

\section{The Main Window at a Glance}
The window is built from a menu bar, a vertical toolbar on the left, several
dockable panels, and a status bar along the bottom.

\begin{table}[h]
\centering
\begin{tabular}{@{}lp{10.5cm}@{}}
\toprule
\textbf{Area} & \textbf{Purpose} \\
\midrule
Menu bar (top) & \ui{Settings} and \ui{View} menus: load/save configs, reset, toggle panels. \\
Left toolbar & Add/remove cameras, display settings, edit ROI, experimental settings. \\
Cameras panel & 2$\times$2 grid of live video with detections and overlays. \\
System Metrics panel & Clock, CPU/RAM/GPU usage, uptime, FPS, latency breakdown. \\
Logs panel & Running list of status messages. \\
Delay Metrics Chart & Capture / inference / display latency over time. \\
System Metrics Chart & CPU / RAM / GPU usage over time. \\
Status bar & Short transient messages about the last action. \\
\bottomrule
\end{tabular}
\end{table}

\subsection{Docking and the View menu}
Every panel is a dock that can be moved, floated, resized, or closed. Drag a
panel's title bar to rearrange it, or drag it out of the window to float it.
The \ui{View} menu lists each panel with a checkbox to show or hide it, plus
\ui{Reset layout} to restore the standard arrangement. Your layout is saved
automatically and restored on the next launch.

\section{The Cameras Panel}
The Cameras panel shows a 2$\times$2 grid of feeds. The four slots are filled in
the order your cameras were added: slots~1--2 on the top row, slots~3--4 on the
bottom row. An empty slot reads \emph{``Missing camera source.''} Drag the
splitter bars between feeds to resize them.

\subsection{What you see on each feed}
Depending on the display settings (\S\ref{sec:display}), each feed can show:
\begin{itemize}[nosep]
  \item \textbf{Bounding boxes} around detected objects, colour-coded:
        green = vehicle, yellow = ambulance, red = firetruck.
  \item \textbf{Labels} (class name and confidence) when annotations are on.
  \item \textbf{An info overlay} (top-left) with camera name, location, and an
        estimated congestion ratio.
  \item \textbf{A traffic-signal indicator} (top-right) showing this lane's
        current red / yellow / green state.
  \item \textbf{The region-of-interest outline} when ROI display is enabled.
\end{itemize}

\section{The Left Toolbar --- Core Actions}
The toolbar holds the five everyday actions. Each is described below.

\subsection{Add source}
Opens a dialog to register a new camera. \emph{Which} dialog appears depends on
the \ui{Use videos} option (\S\ref{sec:experimental}):

\subsubsection{Live camera (RTSP)}
When \ui{Use videos} is off, the \ui{Add Camera Source} form requests:
\begin{itemize}[nosep]
  \item \ui{Camera Name} and \ui{Location} --- free text for identification.
  \item \ui{Max cap} --- the lane's vehicle capacity, used for the congestion
        estimate.
  \item \ui{IP Address} --- validated as a dotted IPv4 address (e.g.\
        \code{192.168.1.64}).
  \item \ui{Username} / \ui{Password} --- camera login (password is masked).
  \item \ui{RTSP Port} --- defaults to \code{554}.
  \item \ui{Channel} --- defaults to \code{101} (main stream of channel~1 on
        Hikvision-style cameras).
  \item \ui{Direction} --- \ui{Vertical} or \ui{Horizontal}; this assigns the
        camera to one of the two traffic lights.
\end{itemize}
The stream URL is assembled as
\code{rtsp://user:pass@ip:port/Streaming/Channels/channel}. All fields are
required; click \ui{Add} to save. The detector restarts automatically with the
new camera.

\subsubsection{Proxy video file}
When \ui{Use videos} is on, the \ui{Add Proxy Video Source} form requests a
\ui{Camera Name}, \ui{Location}, \ui{Max cap}, \ui{Direction}, and a
\ui{Video File} chosen with \ui{Browse} (\code{.mp4}, \code{.avi}, \code{.mkv},
\code{.mov}). This is the recommended way to demo or test the system without
live cameras. Video sources loop continuously.

\subsection{Remove source}
Opens \ui{Remove Camera Source(s)}, showing two checklists side by side:
\ui{Camera Sources} (live) and \ui{Proxy Camera Sources} (video). Tick any
combination across both lists and click \ui{Remove Selected}. The detector
restarts with the remaining sources. \ui{Cancel} discards the selection.

\subsection{Display settings}
\label{sec:display}
Opens \ui{Change Display Settings}, which controls only what is \emph{drawn};
changes apply immediately without restarting the detector. Three groups:

\begin{itemize}[nosep]
  \item \ui{Bounding Box Display} --- show/hide boxes for \ui{Firetrucks},
        \ui{Ambulances}, and \ui{Vehicles} independently. Enable firetruck and
        ambulance boxes when you want to visually confirm emergency detections.
  \item \ui{On-Screen Display (OSD)} --- toggle each overlay element:
        \ui{Name}, \ui{Location}, \ui{Estimated Congestion}, \ui{Annotations}
        (class/confidence labels), \ui{Region of Interests} (the ROI outline),
        and \ui{Traffic Phase} (the red/yellow/green indicator).
  \item \ui{Logs} --- the log verbosity level recorded in the Logs panel:
        \ui{Debug}, \ui{Info}, or \ui{Errors}.
\end{itemize}
Click \ui{Save} to apply, \ui{Cancel} to discard.

\subsection{Edit region (ROI)}
\label{sec:roi}
A region of interest restricts detection to the part of the frame that is
actual road, which improves both accuracy and speed. Click \ui{Edit Region},
pick a camera from the dropdown, and click \ui{OK} to open the ROI editor.

\subsubsection{Using the ROI editor}
The editor shows a snapshot from that camera with a blue quadrilateral overlay:
\begin{itemize}[nosep]
  \item \textbf{Drag a corner handle} (the light-blue squares) to reshape the
        region.
  \item \textbf{Drag inside the polygon} to move the whole region.
  \item The region is clamped to the image bounds.
  \item Click \ui{Save} to store the ROI (kept as normalised coordinates so it
        survives resolution changes), or \ui{Cancel} to discard.
\end{itemize}
Draw the quadrilateral tightly around the lane(s) feeding the intersection.
After saving, enable \ui{Region of Interests} in the display settings to see the
outline on the live feed. Detection thereafter runs only inside this region.

\subsection{Experimental settings}
\label{sec:experimental}
Opens \ui{Change Experimental Settings}. Saving here \emph{restarts} the
detector. Options:
\begin{itemize}[nosep]
  \item \ui{Use videos} --- switch the whole app between live RTSP cameras and
        local video files. This determines which source list is active and which
        \ui{Add source} dialog appears.
  \item \ui{Tile Width} / \ui{Tile Height} --- the resolution each feed is
        processed at. Smaller tiles run faster (useful on a modest GPU such as
        an RTX~2060); larger tiles detect smaller/farther vehicles more
        reliably.
  \item \ui{Model confidence} --- the detection threshold, adjustable from
        \code{0.00} to \code{1.00} via the slider or the spin box (they stay in
        sync). Lower values detect more but risk false positives; the default is
        \code{0.30}.
\end{itemize}

\section{Monitoring Panels}

\subsection{System Metrics}
A compact readout updated once per second: current date and time, \ui{CPU},
\ui{RAM} and \ui{GPU} utilisation bars, \ui{Uptime}, the processing \ui{FPS},
and a latency breakdown --- \ui{Total Delay} and its \ui{Cam}, \ui{Inf} and
\ui{Disp} components (capture, inference, and display time in milliseconds).
FPS is derived as $1/\text{total delay}$.

\subsection{Charts}
Two live charts plot the same data over time so you can spot trends:
\begin{itemize}[nosep]
  \item \ui{Delay Metrics Chart} --- camera, inference, and display latency.
  \item \ui{System Metrics Chart} --- CPU, RAM, and GPU usage.
\end{itemize}
Both keep a rolling window of recent samples.

\subsection{Logs}
A scrolling, read-only history of status messages (most recent first), mirroring
what briefly appears in the status bar --- camera additions/removals, detector
restarts, settings saves, and errors.

\section{The Settings Menu --- Configurations}
Configurations bundle your cameras, display options, experimental options, and
window layout into a single JSON file, so you can save and reload complete
setups (e.g.\ one for the live intersection, one for a video demo).

\begin{itemize}[nosep]
  \item \ui{Load settings} --- import a configuration from a \code{.json} file
        and restart the detector with it. Ready-made files live in
        \code{configs/} (\code{all\_cameras.json} for the live cameras,
        \code{proxy\_cameras.json} / \code{focus\_first\_camera.json} for video
        demos).
  \item \ui{Save settings} --- persist the current state to the app's internal
        store (also happens automatically on most actions).
  \item \ui{Save settings as} --- export the current state to a \code{.json}
        file you choose.
  \item \ui{Reset to default settings} --- restore default display, experimental
        and layout settings (your camera list is kept).
  \item \ui{Reset camera data} --- clear all configured cameras.
  \item \ui{Quit} --- close the application.
\end{itemize}

\section{Signal-Control Behaviour (What the App Decides)}
While running, the app continuously counts vehicles per direction and drives the
two lights. Understanding this helps interpret the on-screen signal indicator.

\begin{itemize}[nosep]
  \item \textbf{Density rule} --- the direction with more vehicles is favoured
        for green; the other is held red.
  \item \textbf{Minimum green} --- a newly granted green is held briefly before
        it can switch, preventing rapid flicker when counts are close.
  \item \textbf{Maximum green} --- a direction cannot hold green indefinitely;
        after the cap, the other direction is given a turn.
  \item \textbf{Yellow transition} --- the losing direction shows yellow before
        going red, then its partner turns green.
  \item \textbf{Emergency preemption} --- a confirmed ambulance or firetruck
        forces its lane green immediately, bypassing the minimum-green hold and
        staying green while the emergency is present. A detection is only
        ``confirmed'' once it clears the emergency confidence threshold for
        several consecutive cycles, so a single noisy frame cannot trigger it.
\end{itemize}

\subsection{Tuning the emergency thresholds}
The two emergency knobs are \emph{not} in the experimental-settings dialog; they
are set in the configuration file's \code{experimental\_settings} block and
applied when the config is loaded:
\begin{lstlisting}
"emergency_confidence": 0.55,   // min confidence for an emergency to count
"emergency_frames": 5           // consecutive cycles required to confirm
\end{lstlisting}
Raise \code{emergency\_confidence} or \code{emergency\_frames} if you see false
preemptions; lower them if genuine emergency vehicles are reacting too slowly.
To watch this live, enable the \ui{Ambulances} and \ui{Firetrucks} bounding
boxes in the display settings.

\section{Typical Workflows}

\subsection{Run a video-file demo (no hardware)}
\begin{enumerate}[nosep]
  \item \ui{Settings $\rightarrow$ Load settings} and pick
        \code{configs/proxy\_cameras.json}, \emph{or}\dots
  \item open \ui{Experimental settings}, tick \ui{Use videos}, \ui{Save}; then
        use \ui{Add source} to add one or more video files.
  \item Set a \ui{Direction} for each so both vertical and horizontal lanes are
        represented.
  \item Optionally draw an ROI per source (\S\ref{sec:roi}).
  \item Watch the feeds, the signal indicators, and the metrics panels.
\end{enumerate}

\subsection{Set up the live intersection}
\begin{enumerate}[nosep]
  \item Ensure \ui{Use videos} is off in \ui{Experimental settings}.
  \item \ui{Add source} for each camera with its IP, credentials, and
        \ui{Direction} (or load \code{configs/all\_cameras.json}).
  \item Draw an ROI for each camera over the relevant lanes.
  \item Tune \ui{Model confidence} and tile size for your GPU.
  \item Confirm the signal indicators react and that the ESP32 lights follow
        (see the Installation Guide's smoke test).
  \item \ui{Settings $\rightarrow$ Save settings as} to store this deployment.
\end{enumerate}

\subsection{Improve performance on a modest GPU}
\begin{enumerate}[nosep]
  \item Lower \ui{Tile Width}/\ui{Height} in \ui{Experimental settings}.
  \item Draw tight ROIs so inference covers less area.
  \item Watch the \ui{Inf Delay} in System Metrics and the Delay Metrics Chart;
        aim for a stable, low inference time and acceptable FPS.
\end{enumerate}

\section{Tips and Troubleshooting}
\begin{itemize}[nosep]
  \item \textbf{A feed says ``Missing camera source''} --- that grid slot has no
        camera; add one, or the stream/file could not be opened (check the IP,
        credentials, network, or file path).
  \item \textbf{Detector seems stuck after changing options} --- adding/removing
        sources and saving experimental settings restart the detector; give it a
        moment to reload the model.
  \item \textbf{Boxes/overlays not showing} --- check the corresponding toggles
        in \ui{Display settings}; these apply instantly.
  \item \textbf{Too many/few detections} --- adjust \ui{Model confidence}.
  \item \textbf{Detections outside the road} --- draw or tighten the ROI.
  \item \textbf{Lost your layout} --- \ui{View $\rightarrow$ Reset layout}.
  \item \textbf{Want a clean slate} --- \ui{Settings $\rightarrow$ Reset to
        default settings} (keeps cameras) or \ui{Reset camera data} (clears
        cameras).
\end{itemize}

\end{document}
